\documentclass[11pt]{article}

\usepackage[margin=1in]{geometry}
\usepackage{amsmath,amssymb}
\usepackage{array}
\usepackage{booktabs}
\usepackage{caption}
\usepackage{enumitem}
\usepackage{float}
\usepackage{graphicx}
\usepackage{hyperref}
\usepackage{multirow}
\usepackage{siunitx}
\usepackage{subcaption}
\usepackage{xcolor}

\hypersetup{
  colorlinks=true,
  linkcolor=blue!55!black,
  citecolor=blue!55!black,
  urlcolor=blue!55!black
}

\captionsetup{font=small, labelfont=bf}
\setlist[itemize]{leftmargin=1.2em, itemsep=0.2em, topsep=0.3em}
\setlist[enumerate]{leftmargin=1.4em, itemsep=0.2em, topsep=0.3em}

\newcommand{\RepoRoot}{../..}
\newcommand{\StepTwoResamplingPlotRoot}{\RepoRoot/surrogate_experiment_results/Step2_resampling/plot_results}
\newcommand{\SPOValidationPlotRoot}{\RepoRoot/surrogate_experiment_results/SPO_validation}
\newcommand{\DecisionAnalysisPlotRoot}{plots}
\newcommand{\howtoread}[1]{%
  \par\smallskip
  \noindent\textbf{How to read.} #1%
}

\title{Weekly Progress: Validation, Resampling, and Decision Analysis}
\author{}
\date{June 3, 2026}

\begin{document}
\maketitle

\begin{quote}
\textbf{Executive summary.}
This week I closed the loop from implementation validation to empirical
variability and decision-level explanation. First, the Step1c SPO+ code was
validated against PyEPO references on both shortest-path and KEP oracles.
Second, subset-resampling boxplots summarize the distribution of 2stage and
SPO+ behavior under fixed graph pools. Third, decision-level analysis on
Step2b d8 suggests that 2stage can remain competitive because many large
edge-level errors are not decision-critical, while SPO+ helps when errors fall
on decision-changing edges.
\end{quote}

\section{Step 1: Implementation Validation with PyEPO Reference}
\label{sec:implementation-validation}

\subsection{Motivation}
The goal of this step is to check whether my implementation contains subtle
bugs. The logic is simple: if my code cannot reproduce PyEPO behavior on a
standard benchmark or on the same KEP oracle, then there may be a bug in the
loss, gradient, optimizer update, solver interface, or metric calculation.

\subsection{Shortest-Path Benchmark}
I first used PyEPO's fixed-data shortest-path setup as an external benchmark.
For the two-stage LR baseline, my implementation fully matches PyEPO across all
fixed-data shortest-path settings:

\begin{table}[h]
\centering
\small
\caption{Shortest-path PyEPO-vs-my comparison.}
\label{tab:shortest-path-validation}
\begin{tabular}{lcc}
\toprule
Quantity & PyEPO LR vs my LR & PyEPO SPO+ vs my SPO+ \\
\midrule
Result files & 24 / 24 & 24 / 24 \\
Rows & 240 / 240 & 240 / 240 \\
Passing settings & 24 / 24 & 19 / 24 \\
True SPO global max diff & $3.12\times 10^{-19}$ & $3.61\times 10^{-4}$ \\
Unambiguous SPO global max diff & $2.07\times 10^{-19}$ & $3.61\times 10^{-4}$ \\
MSE global max diff & $1.09\times 10^{-47}$ & $2.70\times 10^{-3}$ \\
Epoch diff & 0 & 0 \\
\bottomrule
\end{tabular}
\end{table}

\begin{figure}[H]
\centering
\includegraphics[width=\textwidth]{\SPOValidationPlotRoot/PyEPO/images/my-methods-sp-n1000e0.png}
\caption{Shortest-path fixed-data comparison, train size 1000 and noise half-width 0.0.}
\label{fig:shortest-path-n1000-e0}
\howtoread{The blue/orange boxes are PyEPO LR/SPO+, and the green/red boxes are my 2stage LR/SPO+. The near-overlap of PyEPO LR with my 2stage LR shows exact baseline agreement; SPO+ is visually close in this setting, with remaining row-level differences interpreted through the tie-breaking caveat in Table~\ref{tab:shortest-path-validation}.}
\end{figure}

\begin{figure}[H]
\centering
\includegraphics[width=\textwidth]{\SPOValidationPlotRoot/PyEPO/images/my-methods-sp-n1000e5.png}
\caption{Shortest-path fixed-data comparison, train size 1000 and noise half-width 0.5.}
\label{fig:shortest-path-n1000-e5}
\howtoread{This noisier shortest-path setting provides a second visual check. PyEPO LR and my 2stage LR remain aligned across degrees, while PyEPO SPO+ and my SPO+ show the same qualitative pattern; this supports the conclusion that the implementation follows the PyEPO reference up to known shortest-path tie-breaking effects.}
\end{figure}

For SPO+, the loss, gradient, and one-step update match PyEPO when the oracle
solution is shared. The remaining full-CSV differences are explained by
shortest-path tie-breaking: the same cost vector can have multiple optimal
paths, and PyEPO's \texttt{optDataset} can store a different optimal solution
vector across repeated solves even when the optimal objective is identical.
Since SPO+ gradients depend on the solution vector and not only on the objective
value, this can create training divergence without implying an algebraic bug.

\subsection{KEP Reference Validation}
I then moved the validation to the actual KEP oracle. This is more directly
relevant to Step1c because KEP is a reward-maximization problem with
instance-specific graph structure. The KEP validation uses the same
\texttt{CachedHybridKepModel} / Gurobi oracle and compares PyEPO-style
SPOPlus reference computations with my Step1c reward-max SPO+ implementation.

\begin{table}[h]
\centering
\small
\caption{KEP PyEPO-reference vs Step1c validation summary.}
\label{tab:kep-validation}
\begin{tabular}{lcc}
\toprule
Validation phase & Key comparison & Result \\
\midrule
Phase 0 preflight
& reward-max / cost-min sign adapter
& all max diffs $=0$ \\
Phase 1 small setting
& SPO+ forward loss
& max diff $=2.9998\times 10^{-6}$ \\
Phase 1 small setting
& SPO+ gradient and one-step SGD update
& max diff $=0$ \\
Phase 2 full trajectory
& $\theta_1,\theta_2,\|\theta\|$
& max diff $=0$ \\
Phase 2 full trajectory
& train SPO+ loss
& max diff $=3.2975\times 10^{-6}$ \\
Phase 2 full trajectory
& validation SPO+ loss
& max diff $=2.2324\times 10^{-6}$ \\
Phase 2 full trajectory
& train / validation decision gap
& max diff $=0$ \\
\bottomrule
\end{tabular}
\end{table}

\begin{figure}[H]
\centering
\includegraphics[width=\textwidth]{\SPOValidationPlotRoot/kep_vs_pyepo/step2b_bridge_results/plots/step2b_degree_boxplot_train_size=50_seeds=0-9_test_mean_normalized_gap.png}
\caption{KEP Step2b bridge visualization: PyEPO-reference and my implementation on the same degree regimes.}
\label{fig:kep-step2b-bridge-boxplot}
\howtoread{Each degree shows paired boxes for PyEPO-reference LR/SPO+ and my LR/SPO+ on the KEP Step2b bridge setting. The visual check is whether each PyEPO box overlaps its corresponding my-implementation box; the near-identical paired boxes complement the exact max-difference checks in Table~\ref{tab:kep-validation}.}
\end{figure}

\subsection{Conclusion from Step 1}
The implementation check supports the conclusion that the Step1c SPO+
algebra, reward-max sign convention, gradient scaling, optimizer update, KEP
oracle interface, and evaluation path are consistent with an independent
PyEPO-based reference implementation. The remaining shortest-path SPO+ CSV
differences are attributable to oracle tie-breaking among multiple shortest
paths, not to an identified error in the Step1c SPO+ formula.

\section{Step 2: Subset-Resampling Variability}
\label{sec:resampling-variability}

\subsection{Motivation}
The second step is to turn a single-seed observation into distributional
evidence. The current resampling experiment intentionally isolates one source
of randomness: the training subset. The graph pool, label regime, validation
set, heldout set, and training initialization are fixed. Only the subset seed
used to sample 50 training graphs is varied.

Thus, this experiment should be interpreted as:

\begin{quote}
\emph{Train-subset resampling on a fixed graph pool and fixed label regime.}
\end{quote}

It does not yet test graph-generation seed, label-noise seed, or training
initialization seed.

\subsection{Experimental Setup}
The first-pass resampling setting is:

\begin{center}
\begin{tabular}{ll}
\toprule
Item & Setting \\
\midrule
Train size & 50 \\
Randomness varied & subset seed only \\
Graph pool & fixed \\
Training seed & fixed \\
Evaluation & heldout400 \\
Main methods & 2stage selected by validation MSE; SPO+ selected by validation SPO+ loss \\
Primary metric & mean normalized decision gap \\
\bottomrule
\end{tabular}
\end{center}

\subsection{Selected Plots}
The plotting script generated several views of the heldout400 distribution. To
avoid redundant evidence in the report, I keep the two most informative views:
\begin{itemize}
    \item by-block normalized decision gap, which shows the absolute
    performance distribution across regimes;
    \item paired 2stage-minus-SPO+ reduction, which compares the two methods on
    the same subset seeds.
\end{itemize}

\begin{figure}[p]
\centering
\includegraphics[width=\textwidth]{\StepTwoResamplingPlotRoot/phase1_heldout400_by_block_gap_boxplot.pdf}
\caption{Heldout400 normalized decision gap by block under subset resampling.}
\label{fig:step2-heldout400-by-block-gap}
\howtoread{Each box summarizes heldout400 normalized decision gaps across subset seeds for one regime/block and method. Lower boxes indicate better downstream decision quality; wider or higher boxes indicate stronger subset sensitivity.}
\end{figure}

\begin{figure}[p]
\centering
\includegraphics[width=\textwidth]{\StepTwoResamplingPlotRoot/phase1_heldout400_paired_reduction_boxplot.pdf}
\caption{Paired heldout400 improvement: 2stage gap minus SPO+ gap.}
\label{fig:step2-heldout400-paired-reduction}
\howtoread{Each value compares 2stage and SPO+ on the same subset seed. Positive values mean SPO+ has a lower normalized decision gap; the paired view reduces confounding from subset difficulty.}
\end{figure}

\subsection{Interpretation}
These two plots answer whether the SPO+ improvement is stable under different
sampled training subsets from the same fixed graph pool. The by-block gap plot
shows where the absolute gaps are large, and the paired reduction plot is the
main evidence for improvement because it compares 2stage and SPO+ on the same
subset seed, reducing confounding from subset difficulty.

The correct conclusion is limited but useful:
\begin{quote}
SPO+ improvement appears stable under train-subset resampling on the fixed
Step2 graph pool, especially in high-degree misspecification regimes.
\end{quote}

The current experiment should not be described as full robustness over all
randomness. A stronger next-stage robustness test would vary graph-generation
seed, label-noise seed, or training initialization seed.

\clearpage
\section{Step 3: Decision-Level Explanation}
\label{sec:decision-analysis}

\subsection{Research Question and Hypothesis}
After completing the implementation validation and subset-resampling
experiments, the next question is no longer whether the code is correct, but
why the observed decision behavior occurs. In particular, the advisor's
hypothesis was:

\begin{quote}
\emph{The two-stage model may remain competitive because bad edge-level
predictions do not necessarily affect the final KEP solution.}
\end{quote}

Therefore, the goal of this section is to move from aggregate test gaps to
decision-level evidence. The main questions are:

\begin{enumerate}
    \item When 2stage and SPO+ have similar decision quality, are they selecting
    the same KEP solution?
    \item If they select different solutions, are these solutions near-optimal
    alternatives under the true objective?
    \item Are large edge-level prediction errors actually decision-critical, or
    do they often occur on edges that do not affect the selected solution?
    \item When SPO+ improves over 2stage, does it reduce errors on
    decision-changing edges?
\end{enumerate}

The central distinction is between \emph{raw prediction error} and
\emph{decision-critical prediction error}. A large edge prediction error matters
only if it changes the selected KEP solution or appears in the symmetric
difference between the predicted and oracle solutions.

\subsection{Analysis Scope}
This first-pass analysis focuses on the most difficult noiseless polynomial
misspecification setting:

\[
\texttt{Step2b d8}, \qquad n_{\mathrm{train}} = 50,
\qquad \texttt{evaluation} = \texttt{heldout400}.
\]

The two compared methods are:

\[
\text{2stage selected by validation MSE}
\qquad \text{vs.} \qquad
\text{SPO+ selected by validation SPO+ loss}.
\]

This is intentionally a narrow setting. The goal is not to claim a universal
mechanism across all regimes, but to obtain a clean first-pass explanation in a
high-misspecification setting where 2stage and SPO+ differences are meaningful.

\subsection{Analysis Design}
The decision-level analysis has four stages:

\begin{enumerate}
    \item \textbf{Seed selection.} Select representative subset seeds from the
    resampling results: large SPO+ improvement, weak/borderline improvement,
    and easy low-gap cases.

    \item \textbf{Per-graph decision replay.} For each selected seed and each
    heldout graph, replay the 2stage and SPO+ checkpoints and record the
    selected KEP solutions.

    \item \textbf{Edge-level criticality analysis.} For each graph, compute
    edge-level prediction errors and classify each edge according to whether it
    appears in the oracle solution, predicted solution, or their symmetric
    difference.

    \item \textbf{Mechanism extraction.} Use summary plots, correlation
    analysis, near-tie analysis, and case studies to identify why 2stage
    sometimes remains competitive and when SPO+ helps.
\end{enumerate}

The analysis produced the following files:

\begin{center}
\small
\begin{tabular}{lr}
\toprule
Output & Rows / files \\
\midrule
selected case seeds & 9 rows \\
per-graph decision comparison & 7,200 rows \\
edge-level criticality table & 879,642 rows \\
graph-level criticality summary & 7,200 rows \\
case-study index & 9 rows \\
case-study edge tables & 9 tables \\
margin / near-tie analysis & 3,600 rows \\
observed candidate ranking & 10,800 rows \\
decision-critical MSE correlations & 12 rows \\
\bottomrule
\end{tabular}
\end{center}

\subsection{Seed Selection: Why These 9 Seeds}
The selected seeds were chosen from the subset-resampling results, not manually
from individual graphs. The idea is to cover three qualitatively different
situations:

\begin{itemize}
    \item \textbf{Large improvement:} 2stage has relatively high decision gap,
    while SPO+ substantially improves over it.
    \item \textbf{Weak or borderline improvement:} SPO+ does not strongly
    improve over 2stage.
    \item \textbf{Easy low-gap:} both 2stage and SPO+ have low decision gaps.
\end{itemize}

The selected seeds are:

\[
\begin{aligned}
\text{large improvement:} &\quad \{1,25,22\},\\
\text{weak/borderline improvement:} &\quad \{27,21,30\},\\
\text{easy low-gap:} &\quad \{41,16,33\}.
\end{aligned}
\]

This selection makes the later case studies interpretable: we can compare
cases where SPO+ clearly helps, cases where both methods are already good, and
cases where the improvement is weak.

\subsection{Decision Replay: From Aggregate Gap to Per-Graph Decisions}
Aggregate heldout400 gaps are useful, but they do not reveal what decision was
made on each graph. Therefore, for each selected seed, I replayed the trained
2stage and SPO+ checkpoints on every heldout graph and recorded:

\begin{itemize}
    \item the oracle solution $y^\star$ under the true reward
    $w_{\mathrm{true}}$;
    \item the 2stage-selected solution $y_{\mathrm{2stage}}$;
    \item the SPO+-selected solution $y_{\mathrm{SPO+}}$;
    \item true objective values and decision gaps;
    \item solution overlap metrics such as Jaccard similarity with the oracle
    solution.
\end{itemize}

For a method $m$, the decision gap is:

\[
\mathrm{gap}_m
=
w_{\mathrm{true}}^\top y^\star
-
w_{\mathrm{true}}^\top y_m,
\]

and the normalized gap is:

\[
\mathrm{normalized\ gap}_m
=
\frac{
w_{\mathrm{true}}^\top y^\star
-
w_{\mathrm{true}}^\top y_m
}{
\left|w_{\mathrm{true}}^\top y^\star\right| + \epsilon
}.
\]

To ensure that this replay follows the same evaluation path as the original
Step2 heldout400 metrics, the replayed per-graph gaps were checked against each
seed's existing \path{metrics/test_per_graph.csv}. The maximum differences
were:

\[
\max |\Delta \mathrm{gap}| = 5.07\times 10^{-5}, \qquad
\max |\Delta \mathrm{normalized\ gap}| = 1.65\times 10^{-7}.
\]

These differences are within the intended tolerance, so the following analysis
is based on the same decision path as the existing Step2 evaluation.

\subsection{Edge Criticality: From Raw Error to Decision-Changing Error}
The advisor's hypothesis is about edge-level prediction errors. For each edge
$e$, I computed:

\[
\mathrm{err}_e^{(m)} =
\left|
\hat{w}_e^{(m)} - w_{\mathrm{true},e}
\right|,
\]

where $m$ is either 2stage or SPO+. Each edge is then classified using the
following binary indicators:

\begin{center}
\small
\begin{tabular}{ll}
\toprule
Indicator & Meaning \\
\midrule
\texttt{in\_opt} & edge appears in the oracle solution $y^\star$ \\
\texttt{in\_2stage} & edge appears in $y_{\mathrm{2stage}}$ \\
\texttt{in\_spoplus} & edge appears in $y_{\mathrm{SPO+}}$ \\
\texttt{in\_2stage\_symdiff} & edge appears in $y^\star \triangle y_{\mathrm{2stage}}$ \\
\texttt{in\_spoplus\_symdiff} & edge appears in $y^\star \triangle y_{\mathrm{SPO+}}$ \\
\bottomrule
\end{tabular}
\end{center}

The symmetric difference is the key object:

\[
y^\star \triangle y_m
=
\{e: y^\star_e \neq y_{m,e}\}.
\]

Edges in this set are precisely the edges where the predicted solution differs
from the oracle solution. Thus, prediction errors on these edges are more
likely to be decision-critical than prediction errors on edges that are never
selected.

\subsection{First-Pass Aggregate Findings}
Across the 9 selected seeds and 400 heldout graphs, the first-pass summary is:

\begin{table}[H]
\centering
\small
\caption{Decision-level summary on selected Step2b d8 heldout400 cases.}
\label{tab:decision-analysis-summary}
\begin{tabular}{lcc}
\toprule
Metric & 2stage & SPO+ \\
\midrule
Mean normalized decision gap & 0.06087 & \textbf{0.04938} \\
Mean Jaccard with oracle solution & 0.60580 & \textbf{0.63655} \\
Exact oracle-solution rate & 17.81\% & \textbf{23.58\%} \\
Top-10-error edges selected rate & 22.87\% & \textbf{19.05\%} \\
Top-10-error edges in symmetric difference & 9.39\% & \textbf{7.02\%} \\
\midrule
\multicolumn{3}{l}{\textit{Observed 2stage/SPO+ solution relationship}} \\
Same 2stage/SPO+ solution rate & \multicolumn{2}{c}{55.08\%} \\
Mean Jaccard between 2stage and SPO+ & \multicolumn{2}{c}{0.82707} \\
Median Jaccard between 2stage and SPO+ & \multicolumn{2}{c}{1.00000} \\
Any-method near-tie rate, threshold $=0.01$ & \multicolumn{2}{c}{9.17\%} \\
\midrule
\multicolumn{3}{l}{\textit{Correlation with normalized gap, Pearson / Spearman}} \\
All-edge MSE & 0.095 / 0.023 & -0.021 / -0.075 \\
MSE on symmetric-difference edges & 0.390 / 0.528 & 0.272 / 0.372 \\
Top-10-error symmetric-difference rate & \textbf{0.570 / 0.667} & \textbf{0.557 / 0.639} \\
\bottomrule
\end{tabular}
\end{table}

There are three main observations.

First, SPO+ has a lower mean normalized decision gap and higher average solution
overlap with the oracle solution. This confirms that SPO+ is not only improving
a numerical loss but is also selecting edge sets that are closer to the oracle
solution.

Second, the top-10-error symmetric-difference rate is relatively low. For
2stage, only 9.39\% of top-10-error edges enter the symmetric difference with
the oracle solution. For SPO+, this rate is 7.02\%. This suggests that many
large edge-level prediction errors are not decision-changing.

Third, all-edge MSE is weakly correlated with decision gap. In contrast, error
measures restricted to symmetric-difference edges are much more predictive.
This suggests that raw prediction accuracy is not enough to explain decision
quality; what matters is whether errors occur on decision-changing edges.

\subsection{Summary Plots}
The next three figures show the first-pass decision-level diagnostics. They are
kept as separate full-width figures so each diagnostic can be read on its own.

\begin{figure}[p]
\centering
\includegraphics[width=\textwidth]{\DecisionAnalysisPlotRoot/mse_vs_normalized_gap.png}
\caption{Prediction MSE versus normalized decision gap.}
\label{fig:decision-mse-vs-gap}
\howtoread{Each point is a method/graph evaluation from the selected Step2b d8 heldout400 cases. If raw prediction MSE fully explained decision quality, points would align strongly upward; the weak pattern shows that edge error alone is not enough to explain KEP regret.}
\end{figure}

\begin{figure}[p]
\centering
\includegraphics[width=\textwidth]{\DecisionAnalysisPlotRoot/solution_overlap_vs_gap.png}
\caption{Solution overlap with the oracle versus normalized decision gap.}
\label{fig:decision-overlap-vs-gap}
\howtoread{The x-axis measures edge-set overlap with the oracle solution, and the y-axis measures normalized gap. Lower-overlap points with small gaps indicate near-optimal alternative exchanges: the selected solution can differ while remaining close in true objective.}
\end{figure}

\begin{figure}[p]
\centering
\includegraphics[width=\textwidth]{\DecisionAnalysisPlotRoot/high_error_edge_selected_symdiff_rate.png}
\caption{High-error edge selected and symmetric-difference rates.}
\label{fig:decision-high-error-criticality}
\howtoread{The leftmost error-rank bins contain the largest prediction errors. Compare selected-rate and symmetric-difference-rate curves: high-error edges matter most when they enter the selected solution or the predicted-vs-oracle symmetric difference.}
\end{figure}

The MSE-vs-gap plot shows that SPO+ can have higher edge-level MSE but lower
decision gap. The solution-overlap plot shows that lower overlap often
increases gap, but there are also medium-overlap points with small gap,
suggesting near-optimal alternative solutions. The high-error criticality plot
shows that the highest-error edges are selected or enter the symmetric
difference only at limited rates.

\clearpage
\subsection{Case Studies: Three Mechanisms}
To make the aggregate patterns interpretable, I selected three types of graph
case studies.

\subsubsection*{Case A: Bad Prediction but Irrelevant}

\begin{figure}[H]
\centering
\small
\fbox{
\begin{minipage}{0.97\textwidth}
\textbf{Case A: Bad prediction but irrelevant.}

\vspace{0.5em}
\begin{tabular}{lcccc}
\toprule
Seed & Graph & 2stage gap & All-edge MSE & Top-10-error symdiff rate \\
\midrule
41 & \texttt{G-234.json}  & 0 & 19.6328 & 0.0 \\
41 & \texttt{G-994.json}  & 0 & 18.1776 & 0.0 \\
25 & \texttt{G-1516.json} & 0 & 17.7713 & 0.0 \\
\bottomrule
\end{tabular}

\vspace{0.8em}
These graphs have large all-edge prediction error under the 2stage model, but
the selected solution is still optimal. The top-error edges do not enter the
optimal-vs-predicted symmetric difference. Therefore, the errors are large in
prediction space but irrelevant in decision space.
\end{minipage}
}
\caption{Large prediction errors can be non-critical.}
\label{fig:case-a}
\howtoread{These rows have high all-edge MSE but zero 2stage decision gap. The key column is the top-10-error symmetric-difference rate: it is zero, so the largest prediction errors do not affect the decision-changing edge set.}
\end{figure}

This case directly supports the hypothesis that bad edge estimates do not
necessarily affect the final solution.

\subsubsection*{Case B: Different Solution but Near-Optimal}

\begin{figure}[H]
\centering
\small
\fbox{
\begin{minipage}{0.97\textwidth}
\textbf{Case B: Different solution but near-optimal.}

\vspace{0.5em}
\begin{tabular}{lcccc}
\toprule
Seed & Graph & 2stage gap & Jaccard with oracle & Same as oracle? \\
\midrule
1 & \texttt{G-696.json}  & 0.000644 & 0.6818 & No \\
1 & \texttt{G-1372.json} & 0.001174 & 0.6087 & No \\
1 & \texttt{G-607.json}  & 0.001175 & 0.7143 & No \\
\bottomrule
\end{tabular}

\vspace{0.8em}
Here, the 2stage solution is not identical to the oracle solution, and the
Jaccard overlap is only moderate. However, the true objective gap is nearly
zero. This suggests that the graph contains near-optimal alternative exchanges.
\end{minipage}
}
\caption{Different selected solutions can be near-optimal alternatives.}
\label{fig:case-b}
\howtoread{These rows are not exact oracle solutions and have only medium Jaccard overlap, but their normalized decision gaps are close to zero. This supports the near-tie story: solution identity can change while true objective remains almost unchanged.}
\end{figure}

This explains another way in which 2stage can remain competitive: it may select
a different exchange, but that exchange is nearly equivalent under the true
objective.

\subsubsection*{Case C: SPO+ Fixes Decision-Critical Errors}

\begin{figure}[H]
\centering
\small
\fbox{
\begin{minipage}{0.97\textwidth}
\textbf{Case C: SPO+ fixes 2stage decision-critical errors.}

\vspace{0.5em}
\begin{tabular}{lccccc}
\toprule
Seed & Graph & 2stage gap & SPO+ gap
& 2stage top-10 symdiff & SPO+ top-10 symdiff \\
\midrule
1  & \texttt{G-392.json}  & 0.313849 & 0        & 0.3 & 0.0 \\
30 & \texttt{G-1560.json} & 0.292305 & 0.011964 & 0.3 & 0.0 \\
1  & \texttt{G-39.json}   & 0.316226 & 0.080459 & 0.2 & 0.1 \\
\bottomrule
\end{tabular}

\vspace{0.8em}
In these graphs, the 2stage model makes errors that affect decision-changing
edges. SPO+ substantially reduces the decision gap and also reduces the
involvement of high-error edges in the symmetric difference.
\end{minipage}
}
\caption{SPO+ helps when prediction errors are decision-critical.}
\label{fig:case-c}
\howtoread{Compare the 2stage and SPO+ gap columns together with the top-10-error symmetric-difference rates. SPO+ reduces the gap precisely in cases where 2stage's large errors enter decision-changing edges.}
\end{figure}

This case explains when SPO+ is useful: it helps most when prediction errors
fall on the edges that change the selected KEP solution.

\subsection{Best-vs-Second Check for the Selected Cases}
The previous case studies compare the model-selected rank-1 solution against
the oracle solution. To check whether these examples also support the
near-optimal-alternative explanation, I computed the best and second-best
distinct edge-selection solutions under each model's predicted weights.

There are two different margins here and they should not be mixed:

\begin{itemize}
    \item the \textbf{oracle gap}
    $w_{\mathrm{true}}^\top y^\star - w_{\mathrm{true}}^\top y$, which says
    how far a rank-1 or rank-2 solution is from the true oracle;
    \item the \textbf{predicted best-second margin}
    $\hat w^\top y_{\mathrm{rank1}}-\hat w^\top y_{\mathrm{rank2}}$, which
    says how much the model itself prefers rank 1 over rank 2.
\end{itemize}

The main quantity for the story is the oracle gap. The predicted margin is
only an auxiliary diagnostic: a small predicted margin means the model itself
sees the top two solutions as close.

\begin{figure}[H]
\centering
\includegraphics[width=\textwidth]{\DecisionAnalysisPlotRoot/case_best_second_gap_summary.png}
\caption{Rank-1 and rank-2 normalized oracle gaps for the selected Case A/B/C graphs.}
\label{fig:case-best-second-gap}
\howtoread{Each panel shows one method. Rank 1 is the model's best predicted KEP solution and rank 2 is the best distinct edge-selection solution after excluding rank 1. The dashed line marks a 5\% normalized oracle gap.}
\end{figure}

\begin{table}[H]
\centering
\small
\caption{Case-level rank-1 versus rank-2 normalized oracle gaps.}
\label{tab:case-best-second-gap}
\resizebox{\textwidth}{!}{%
\begin{tabular}{lllrrrr}
\toprule
Case & Graph & Method & Rank-1 gap & Rank-2 gap & $\Delta$ gap & Pred. margin \\
\midrule
A1 & G-234 & 2stage & 0.0000 & 0.0074 & 0.0074 & 3.026 \\
A1 & G-234 & SPO+ & 0.0000 & 0.0074 & 0.0074 & 4.726 \\
A2 & G-994 & 2stage & 0.0000 & 0.0210 & 0.0210 & 4.135 \\
A2 & G-994 & SPO+ & 0.0000 & 0.0927 & 0.0927 & 3.102 \\
A3 & G-1516 & 2stage & 0.0000 & 0.1088 & 0.1088 & 0.018 \\
A3 & G-1516 & SPO+ & 0.0000 & 0.0033 & 0.0033 & 0.323 \\
B1 & G-696 & 2stage & 0.0006 & 0.0102 & 0.0096 & 5.751 \\
B1 & G-696 & SPO+ & 0.0006 & 0.0375 & 0.0369 & 3.605 \\
B2 & G-1372 & 2stage & 0.0012 & 0.0000 & -0.0012 & 1.096 \\
B2 & G-1372 & SPO+ & 0.0012 & 0.0000 & -0.0012 & 1.752 \\
B3 & G-607 & 2stage & 0.0012 & 0.0404 & 0.0392 & 0.242 \\
B3 & G-607 & SPO+ & 0.0012 & 0.0404 & 0.0392 & 1.051 \\
C1 & G-392 & 2stage & 0.3138 & 0.3194 & 0.0055 & 0.498 \\
C1 & G-392 & SPO+ & 0.0000 & 0.0253 & 0.0253 & 2.009 \\
C2 & G-1560 & 2stage & 0.2923 & 0.0120 & -0.2803 & 2.493 \\
C2 & G-1560 & SPO+ & 0.0120 & 0.2900 & 0.2780 & 1.543 \\
C3 & G-39 & 2stage & 0.3162 & 0.3082 & -0.0081 & 1.096 \\
C3 & G-39 & SPO+ & 0.0805 & 0.0850 & 0.0046 & 1.274 \\
\bottomrule
\end{tabular}

}
\end{table}

The selected cases now separate into the intended mechanisms. In Case A, rank
1 is exactly oracle for both methods. Rank 2 is often still near-oracle, but
not always; either way, the large edge prediction errors did not change the
top decision. In Case B, rank 1 is not oracle but has near-zero gap, and all
rank-2 solutions stay within 5\% oracle gap. This is the strongest case-level
evidence for a flat near-optimal plateau. In Case C, 2stage rank 1 is bad in
all three examples, but rank 2 distinguishes failure modes: \texttt{G-1560}
has a near-oracle 2stage rank-2 solution, consistent with a local ranking flip,
whereas \texttt{G-392} and \texttt{G-39} keep high 2stage rank-2 gaps,
suggesting that the top predicted region is also poor. SPO+ moves rank 1 back
toward the near-oracle region, but its rank 2 can be much worse, showing that
top-1 ordering matters in some decision-critical cases.

\subsection{Near-Tie and Margin Analysis}
The case studies suggest that some different selected solutions may be
near-optimal alternatives. To quantify this, I compared the observed candidate
solutions:

\[
y^\star,\qquad y_{\mathrm{2stage}},\qquad y_{\mathrm{SPO+}}.
\]

This is not an enumeration of all feasible KEP solutions; it only ranks the
solutions actually observed in the replay. With a near-tie threshold of 0.01 in
normalized gap, the first-pass summary is:

\begin{center}
\small
\begin{tabular}{lr}
\toprule
Quantity & Value \\
\midrule
Paired graph rows & 3,600 \\
Observed candidate rows & 10,800 \\
Mean $| \mathrm{obj}_{2stage} - \mathrm{obj}_{SPO+} |$ & 4.3366 \\
Median $| \mathrm{obj}_{2stage} - \mathrm{obj}_{SPO+} |$ & 0 \\
Mean Jaccard$(2stage, SPO+)$ & 0.8271 \\
Median Jaccard$(2stage, SPO+)$ & 1.0000 \\
Same 2stage/SPO+ solution rate & 55.08\% \\
2stage different-solution near-tie rate & 6.08\% \\
SPO+ different-solution near-tie rate & 6.58\% \\
Any-method near-tie rate & 9.17\% \\
\bottomrule
\end{tabular}
\end{center}

This supports the near-optimal-alternative explanation. In more than half of
the observed graph-seed pairs, 2stage and SPO+ select the same solution. When
they differ, a nontrivial subset of cases still has normalized gap below 1\%,
meaning that different solution identities can correspond to very similar true
objectives.

\subsection{Decision-Critical MSE Correlations}
Finally, I tested whether different prediction-error summaries explain
normalized decision gap. The main comparison is between all-edge MSE and
decision-focused error summaries such as MSE on symmetric-difference edges.

\begin{table}[H]
\centering
\small
\caption{Correlation between prediction-error summaries and normalized decision gap.}
\label{tab:critical-mse-correlations}
\begin{tabular}{p{0.38\linewidth}cc}
\toprule
Predictor & 2stage Pearson / Spearman & SPO+ Pearson / Spearman \\
\midrule
All-edge MSE & 0.095 / 0.023 & -0.021 / -0.075 \\
MSE on selected edges & -0.260 / -0.237 & -0.112 / -0.066 \\
MSE on symmetric-difference edges & 0.390 / 0.528 & 0.272 / 0.372 \\
Top-10-error symmetric-difference rate & \textbf{0.570 / 0.667} & \textbf{0.557 / 0.639} \\
\bottomrule
\end{tabular}
\end{table}

The strongest predictor for both methods is the rate at which top-error edges
enter the symmetric difference. This provides the most direct evidence for the
decision-critical-error hypothesis: prediction error matters primarily when it
changes the selected solution.

\subsection{Answer to the Advisor's Hypothesis and Limitations}
The first-pass evidence supports the hypothesis. The reason 2stage can remain
competitive is not that it predicts every edge accurately. Instead, many large
edge-level errors are non-critical because they occur on edges that are not
selected or do not enter the symmetric difference. In addition, KEP graphs may
contain near-optimal alternative solutions, so two different selected edge sets
can have very similar true objectives.

SPO+ appears most helpful when the 2stage prediction errors fall on
decision-changing edges. In those cases, SPO+ reduces both the decision gap and
the involvement of high-error edges in the symmetric difference.

This is still a first-pass analysis. The current limitations are:

\begin{itemize}
    \item only train size 50 and heldout400 are used;
    \item only 9 representative subset seeds are selected;
    \item \texttt{Step2c} multiplicative-noise regimes are not yet analyzed;
    \item the near-tie analysis uses only observed candidates
    $y^\star$, $y_{\mathrm{2stage}}$, and $y_{\mathrm{SPO+}}$;
    \item the best-vs-second analysis only inspects the top two predicted
    distinct edge-selection solutions, not the full feasible solution set;
    \item adversarial SPO+ solutions $y_{\mathrm{adv}}$ are not yet included.
\end{itemize}

The next step is to extend this same analysis to \texttt{Step2c d8 eps050}, add
FY, and include $y_{\mathrm{adv}}$ to inspect how SPO+ gradients push the
predictions away from bad decision-changing solutions.

\section{Overall Conclusions and Next Steps}
\label{sec:weekly-next-steps}

\subsection{Conclusions}
This week's work forms a closed loop:

\begin{enumerate}
    \item \textbf{Code validation.} The Step1c SPO+ implementation aligns with
    PyEPO reference computations on both shortest-path and KEP oracles. The
    remaining shortest-path full-training SPO+ discrepancy is explained by
    oracle tie-breaking among multiple shortest paths.
    
    \item \textbf{Distributional evidence.} The subset-resampling box plots show
    how 2stage and SPO+ vary across different 50-graph training subsets on a
    fixed graph pool.
    
    \item \textbf{Mechanistic explanation.} The decision-level analysis suggests
    why 2stage can remain competitive: many prediction errors are not
    decision-critical, and some different selected exchanges are near-optimal.
    SPO+ improves most when errors occur on decision-changing edges.
\end{enumerate}

\begin{table}[h]
\centering
\small
\caption{Safe claims supported by the current evidence.}
\label{tab:safe-claims}
\begin{tabular}{p{0.29\linewidth}p{0.39\linewidth}p{0.22\linewidth}}
\toprule
Claim & Evidence & Limitation \\
\midrule
Step1c SPO+ implementation is reliable.
& PyEPO-reference validation on shortest-path and KEP; KEP Phase 2 trajectory equivalence.
& Shortest-path full CSV SPO+ is affected by oracle tie-breaking. \\
\midrule
SPO+ improves over 2stage in high-degree subset-resampling settings.
& Heldout400 boxplots and paired reductions on fixed graph pools.
& Only subset seed varied; not graph seed or label seed. \\
\midrule
2stage remains competitive because many errors are non-critical.
& Case A/B/C, high-error symdiff rates, and decision-critical MSE correlations.
& First-pass evidence on selected Step2b d8 seeds only. \\
\bottomrule
\end{tabular}
\end{table}

\end{document}
